\chapter{Changelog summary}\label{ch:changelog}

TandemApp shipped inside the engine's releases until the two were split (ENG-413). These are the paragraphs on the Chat tab, the dashboard and the proxy from the engine's changelog chapter; the rest of each release is in \enginemanual{Changelog summary}.

\section*{0.10.0, Nemotron 3.5 Bolt}

\paragraph{Chat tab.} Each chat's sandbox has 1~GiB of memory, and a call the kernel kills at the limit says so. An attachment is never inlined past the window: every attachment is saved under the chat's \code{/workspace/uploads/}, and a text file past a quarter of the free window goes as one line with its path. A tool call no longer waits without end for its sandbox. \code{POST /v1/cache/release} gives back memory the engine holds idle.

\section*{0.9.1, web search, media and the loop guard}

\paragraph{Chat tab.} A built-in \code{web_search} tool with DuckDuckGo and SearXNG providers, set in the chat settings. Video, audio and PDF files open in a player or the browser's PDF viewer, in the Files pane and the \code{present} card; any other binary file shows its size and a Download button. chatd streams a workspace file with byte ranges (206, 416, ETag) and no 2~MiB cap.

\section*{0.8.2, community weights and permission modes}

\paragraph{Permission modes.} Each chat in the Chat tab runs in one of three modes, and Auto is the new default: a call that stays inside the chat's sandbox runs without a prompt, and one that leaves it still asks. Deny rules and plan mode's read-only list hold in every mode. A tool card that ran without a prompt carries a badge, and every such call is logged as \code{tool.auto_approved}. \enquote{Always allow} on a chained bash command now saves a rule for every part not yet allowed; before, one chat saved \code{cd} 11 times and still got the prompt. Chapter~\ref{ch:chat}.

\section*{0.8.0, Qwen3.8-Flash-Next and the Chat tab}

\paragraph{The Chat tab.} A chat in the dashboard backed by chatd, its own service behind the engine's \code{/v1/chat-app}: runs that outlive the browser, tools in a sandbox per chat, the workspace panes, MCP servers, plan mode, the PWA. The Context panel reads the usage frame's split by part and the cache peek; a chat on the NVMe is read back while the engine is idle (\code{/v1/cache/prewarm}). Found by the live checks and fixed: a restart's restore that faulted the GPU on Flash-Next, the side lane on Flash-Next, its missing model card, a side slot's chat that came back cold after a restart, a peek that answered 409 right after a turn, a prewarm that held a request for its whole read (it now stops between blocks), a chat waiting for a permission prompt that held one of chatd's two run slots, and an iPhone test, where five screenshots counted by their base64 filled the window (319,136 tokens estimated against 16,980 real). The Chat tab carries a Beta pill in this release. On a phone the bottom bar has four slots with the rest behind More, and the composer shows every control.

\paragraph{Chat tab (beta): known limitations.} The Chat tab is in beta in this release, and a Beta pill beside it says so. We know of these limits and have not fixed them yet:
\begin{itemize}
  \item The agent loop does not yet work the way Claude Code does (epic VIS-E10). Its prompts are not kept prefix-stable from step to step (SRV-375). Old tool results and old thinking stay in the context until a compaction (SRV-365). A compaction writes a general summary, not the task's state (SRV-373). Every sandbox command asks for permission; there is no mode that runs them without asking (SRV-387). Nothing notices when the model repeats the same failing approach with different calls (SRV-367). \code{render_page} passes on Chromium's console noise (SRV-369). Long tool output has no ``see more'' paging (SRV-371). Searches run in the main context, with no sub-agents (SRV-381).
  \item The engine reads a short Flash-Next chat again on its next turn when the answer ended inside a verified block (ENG-235).
  \item A request that arrives during a prewarm still waits 0.5 to 1 s (SRV-349).
  \item When memory reaches its ceiling, the engine drops the parked chat without writing it to the NVMe, so its next turn starts cold.
  \item Images only: there is no video input (VIS-97).
  \item The 5.05 Flash-Next pack does not load yet.
  \item Parallel decoding does not raise throughput over one request at a time: three requests together took 27.18 s against 27.03 s one after another, though the later ones start sooner.
  \item The phone layout's keyboard behaviour is not verified in Safari on an iPhone.
\end{itemize}

\section*{What 0.4.0 brought}

\paragraph{The dashboard.} The Conductor view with all its screens, the Parallel page, the hosts and the memory plan in the System view, the Backends and Setup views, filters by backend and user, and a phone layout for every view. TandemLLM's own calls are left out of Usage unless asked, and the Live panel's cached tokens never read above the prompt. It asks for no password: the admin routes are open until an admin password is set, so the engine belongs on a trusted network (\enginemanual{Access and sign-in}).

\paragraph{The proxy.} \code{tandem-dash}: the dashboard in front of any OpenAI-compatible engine or cloud API, six backend kinds, routing by model id, fallbacks, one config database overlaid by \code{TANDEM_*} variables, encrypted keys, per-user keys, cost at each model's price; vLLM's and SGLang's own metrics, the plugins inside them and \code{tandem-dash launch}; the host collector; the conductor in front of any backend, with the memory writer in the backend's idle time.
